\documentclass[11pt,a4paper]{article}

\usepackage[T1]{fontenc}
\usepackage[utf8]{inputenc}
\usepackage{lmodern}

\usepackage[a4paper,margin=1in]{geometry}
\usepackage{setspace}
\setstretch{1.12}

\setlength{\parindent}{15pt}
\setlength{\parskip}{0pt}

\usepackage{lmodern}
\usepackage{amsmath,amssymb,amsfonts}
\usepackage{booktabs}
\usepackage{longtable}
\usepackage{array}
\usepackage{tabularx}
\usepackage{enumitem}
\usepackage{hyperref}
\usepackage{setspace}
\usepackage{ragged2e}

\hypersetup{
	colorlinks=true,
	linkcolor=black,
	citecolor=black,
	urlcolor=blue,
	pdftitle={Dependency Tables and Node Registry for the GTCS},
	pdfauthor={Kostiantyn Osmolovskyi}
}

\setlist[itemize]{topsep=3pt,itemsep=2pt,parsep=0pt}
\setlist[enumerate]{topsep=3pt,itemsep=2pt,parsep=0pt}


\title{Dependency Tables and Node Registry\\ for the General Theory of Cognitive Structuring}
	
\author{
	Kostiantyn Osmolovskyi\\
	\small Independent Researcher, Ukraine\\
	\small ORCID: 0009-0006-3144-7237\\
	\small \texttt{constantinosmol@gmail.com}
}

\date{}

% ---------------------------------------------------------------------------
% Part of a series: General Theory of Cognitive Structuring (GTCS)
% Autor: Kostiantyn Osmolovskyi (ORCID: 0009-0006-3144-7237)
% License: Creative Commons Attribution 4.0 International (CC BY 4.0)
% Repository: https://github.com/k-osmolovskyi/k-osmolovskyi.github.io
% DOI: https://doi.org/10.5281/zenodo.19701848
% ---------------------------------------------------------------------------

\begin{document}
	\maketitle

\begin{center}
	\small
	Technical Note No. 2026-4\\
	Version 1.0
\end{center}

\vspace{0.5em}
	
	\maketitle
	
	\section{Purpose of this document}
	
	This document provides a structured dependency companion for the paper series associated with the
	\emph{General Theory of Cognitive Structuring}. Its purpose is to make explicit the current working
	registry of dependency nodes and the paper-level dependency structure used to support external
	verification of the series.
	
	The document is intended to accompany, rather than replace, the public
	\emph{Acyclicity Statement and Dependency Criteria} note. That statement specifies the criteria
	under which the series may be treated as a directed acyclic structure. The present document adds
	the concrete registry and dependency tables that make such evaluation traceable at the level of
	papers, definitions, operators, constructions, and proposition-level results.
	
	This document is not a proof compendium and is not intended to provide a complete bibliography
	of all local cross-references. It is a dependency-oriented reconstruction.
	
	\section{Scope and status}
	
	The present registry is versioned as a working public companion.
	
	It is designed to serve four functions:
	\begin{itemize}[leftmargin=2em]
		\item to identify the main dependency nodes used across the series;
		\item to distinguish first appearance from primary formalization;
		\item to separate stable nodes from boundary cases requiring further textual refinement;
		\item to support later graph construction at the node, paper, and layer levels.
	\end{itemize}
	
	The document uses the following status labels:
	\begin{itemize}[leftmargin=2em]
		\item \textbf{S} --- stable;
		\item \textbf{V} --- needs textual verification or boundary refinement;
		\item \textbf{G} --- graph-organizing node, useful for dependency structure but not always a
		basic theorem-grade unit.
	\end{itemize}
	
	\section{Layer scheme}
	
	For dependency purposes, the current series is organized into five broad layers:
	
	\begin{itemize}[leftmargin=2em]
		\item \textbf{L1} --- Entry distinctions;
		\item \textbf{L2} --- Core formal geometry and admissibility;
		\item \textbf{L3} --- Structural architecture and invariants;
		\item \textbf{L4} --- Access and representation;
		\item \textbf{L5} --- Trajectory, inter-system extension, and synthesis.
	\end{itemize}
	
	These layers are organizational rather than ontological. Their purpose is to make the unfolding
	order of the series explicit and readable.
	
	\section{Node registry table v1}
	
	\begin{longtable}{>{\RaggedRight\arraybackslash}p{0.03\textwidth}
			>{\RaggedRight\arraybackslash}p{0.19\textwidth}
			>{\RaggedRight\arraybackslash}p{0.05\textwidth}
			>{\RaggedRight\arraybackslash}p{0.05\textwidth}
			>{\RaggedRight\arraybackslash}p{0.18\textwidth}
			>{\RaggedRight\arraybackslash}p{0.18\textwidth}
			>{\RaggedRight\arraybackslash}p{0.18\textwidth}
			>{\RaggedRight\arraybackslash}p{0.04\textwidth}}
		\toprule
		\textbf{ID} & \textbf{Name} & \textbf{Type} & \textbf{Layer} & \textbf{First appearance} & \textbf{Primary formalization} & \textbf{Direct dependencies} & \textbf{Status} \\
		\midrule
		\endfirsthead
		\toprule
		\textbf{ID} & \textbf{Name} & \textbf{Type} & \textbf{Layer} & \textbf{First appearance} & \textbf{Primary formalization} & \textbf{Direct dependencies} & \textbf{Status} \\
		\midrule
		\endhead
		
		N001 & Admissibility of structural updating & D & L1 & \emph{Structural Updating and the Limits of Cognitive Change} & same paper & --- & S \\
		
		N002 & Structural updating & D & L1 & \emph{Structural Updating and the Limits of Cognitive Change} & \emph{Structural Admissibility in Cognitive Systems} & N001 & S \\
		
		N003 & Processing access / processing without structural updating & D & L1 & \emph{Structural Updating and the Limits of Cognitive Change} & same paper & --- & V \\
		
		N004 & Distinction between processing and structural updating & P & L1 & \emph{Structural Updating and the Limits of Cognitive Change} & same paper & N002, N003 & S \\
		
		N010 & Configuration space $X$ & C & L2 & \emph{Coherence Evaluation in Cognitive Systems} & same paper & --- & S \\
		
		N011 & Current configuration $x_t$ & C & L2 & \emph{Coherence Evaluation in Cognitive Systems} & same paper & N010 & S \\
		
		N012 & Stability / admissible region $U$, later $U_X(I_t)$ & C & L2 & \emph{Coherence Evaluation in Cognitive Systems} & \emph{Invariants in Cognitive Architectures} & N010, N030/N031 & V \\
		
		N013 & Coherence & D & L2 & \emph{Coherence Evaluation in Cognitive Systems} & same paper & N011, N012 & S \\
		
		N014 & Coherence measure $C_I(x_t)$ & O & L2 & \emph{Coherence Evaluation in Cognitive Systems} & same paper & N011, N012, N013 & S \\
		
		N015 & Coherence evaluation & C & L2 & \emph{Coherence Evaluation in Cognitive Systems} & same paper & N013, N014 & S \\
		
		N016 & Structural tension & D & L2 & \emph{Coherence Evaluation in Cognitive Systems} as $T_t = K_t + C_t$; developed later & \emph{Structural Admissibility in Cognitive Systems} & N017, N018 & V \\
		
		N017 & Coordination cost $C_t$ & D & L2 & \emph{Coherence Evaluation in Cognitive Systems} & \emph{Structural Admissibility in Cognitive Systems} & --- & V \\
		
		N018 & Incompatibility load $K_t$ & C & L2 & \emph{Coherence Evaluation in Cognitive Systems} / \emph{Structural Admissibility in Cognitive Systems} & \emph{Invariants in Cognitive Architectures} & N034, N036 & V \\
		
		N019 & Tension decomposition $T_t = K_t + C_t$ & C & L2 & \emph{Coherence Evaluation in Cognitive Systems} & \emph{Structural Admissibility in Cognitive Systems} & N016, N017, N018 & S \\
		
		N020 & Resource state $R_t$ & C & L2 & \emph{Structural Admissibility in Cognitive Systems} & same paper & --- & S \\
		
		N021 & Rigidity / mean rigidity $\bar h_t$ & C & L2 & \emph{Structural Admissibility in Cognitive Systems} & same paper & --- & S \\
		
		N022 & Compensability threshold $\theta(R_t,\bar h_t)$ & O & L2 & \emph{Overload Formation in Cognitive Processing} & \emph{Structural Admissibility in Cognitive Systems} & N020, N021 & S \\
		
		N023 & Instantaneous overload $u_t$ & C & L2 & \emph{Overload Formation in Cognitive Processing} & same paper & N016, N022 & S \\
		
		N024 & Overload memory $L_t$ & C & L2 & \emph{Overload Formation in Cognitive Processing} & same paper & N023 & S \\
		
		N025 & Trajectory dependence & D & L2/L5 & \emph{Overload Formation in Cognitive Processing} & \emph{Trajectory-Dependent Regulation in Cognitive Systems} & N024 & S \\
		
		N026 & Admissibility operator $A_t$ & O & L2 & \emph{Structural Admissibility in Cognitive Systems} & same paper & N016, N020, N021, N024 & S \\
		
		N027 & Stability region in regulatory phase space $U(S_t)$ / $U(DS)$ & C & L2 & \emph{Structural Admissibility in Cognitive Systems} & same paper & N026 & S \\
		
		N028 & Critical boundary / critical tension function $T_{\mathrm{crit}}(L)$ & C & L2 & \emph{Structural Admissibility in Cognitive Systems} & \emph{Trajectory-Dependent Regulation in Cognitive Systems} & N026, N027 & S \\
		
		N029 & Architectural level & D & L2 & \emph{Structural Admissibility in Cognitive Systems} & same paper & N027 & S \\
		
		N029a & Level transition & D/P & L2/L3 & \emph{Structural Admissibility in Cognitive Systems} & \emph{Structural Compression in Cognitive Architectures} & N029, N027, N039 & S \\
		
		N030 & Invariant & D & L3 & implicit earlier; explicit only later & \emph{Invariants in Cognitive Architectures} & --- & S \\
		
		N031 & Invariant structure $I_t$ & C & L3 & implicit earlier; explicit only later & \emph{Invariants in Cognitive Architectures} & N030 & S \\
		
		N032 & Interaction graph $G_t$ & C & L3 & \emph{Structural Admissibility in Cognitive Systems} state description & \emph{Invariants in Cognitive Architectures} & N031 & S \\
		
		N033 & Local admissibility profile $a_i(x)$ & O & L3 & \emph{Structural Admissibility in Cognitive Systems} soft compatibility & \emph{Invariants in Cognitive Architectures} & N030, N010 & S \\
		
		N034 & Pairwise incompatibility $c_{ij}(x)$ & O & L3 & \emph{Structural Admissibility in Cognitive Systems} & \emph{Invariants in Cognitive Architectures} & N032, N033 & S \\
		
		N035 & Structural conflict & D & L3 & implicit earlier & \emph{Invariants in Cognitive Architectures} & N034 & V \\
		
		N036 & Global conflict aggregation $K_t$ & C & L3 & \emph{Structural Admissibility in Cognitive Systems} & \emph{Invariants in Cognitive Architectures} & N034, N035 & S \\
		
		N037 & Invariant-induced stability region $U_X(I_t)$ & C & L3 & implicit in coherence paper as $U$ & \emph{Invariants in Cognitive Architectures} & N031, N033 & S \\
		
		N038 & Identity of invariant-constrained organization & M & L3/L5 & \emph{Identity as a Regulatory Attractor} & same paper & N031, N027 & G \\
		
		N039 & Compression operation $I_{t+1} = (I_t \setminus S) \cup \{F(S)\}$ & O & L3 & implicit earlier & \emph{Structural Compression in Cognitive Architectures} & N031, N026 & S \\
		
		N040 & Compressed invariant $F(S)$ & C & L3 & \emph{Structural Compression in Cognitive Architectures} & same paper & N039 & S \\
		
		N041 & Compression-induced geometry change & P & L3 & \emph{Structural Compression in Cognitive Architectures} & same paper & N039, N027 & S \\
		
		N042 & Compression-induced level transition criterion & P & L3 & \emph{Structural Compression in Cognitive Architectures} & same paper & N041, N029a & S \\
		
		N050 & Pre-symbolic admissibility & D & L4 & implicit layered version in \emph{General Theory} & \emph{Pre-Symbolic Admissibility in Cognitive Systems} & N026, N031, N057 & S \\
		
		N051 & Pre-symbolic admissibility operator $\Pi_t^{ps}$ & O & L4 & \emph{Pre-Symbolic Admissibility in Cognitive Systems} & same paper & N050 & S \\
		
		N052 & Pre-symbolic signal domain $Z_t$ & C & L4 & \emph{Pre-Symbolic Admissibility in Cognitive Systems} & same paper & --- & S \\
		
		N053 & Admissible discrepancy domain $D_t = Z_t^{adm}$ & C & L4 & \emph{General Theory} in synthetic form; explicit formalization later & \emph{Pre-Symbolic Admissibility in Cognitive Systems} & N051, N052 & S \\
		
		N054 & Regulatory accessibility & D & L4 & implied by pre-symbolic filtering & \emph{Restricted Accessibility of Coherence in Cognitive Systems} & N050, N053 & S \\
		
		N055 & Restricted accessibility of coherence & P & L4 & \emph{Restricted Accessibility of Coherence in Cognitive Systems} & same paper & N013, N053, N054 & S \\
		
		N056 & Non-injectivity of regulatory accessibility & P & L4 & \emph{Restricted Accessibility of Coherence in Cognitive Systems} & same paper & N053, N055 & S \\
		
		N057 & Coherence representation & D & L4 & \emph{Emergence of Coherence Representation} & same paper & N013, N024, N054 & S \\
		
		N058 & Expressed regulatory representation $\widehat{C}_t$ / $Cb_t$ & R & L4 & \emph{Emergence of Coherence Representation} & same paper & N057, N060 & S \\
		
		N059 & Extended regulatory representation $\widetilde{C}_t^A$ / related derived variable & R & L4 & later representational papers & \emph{Coherence Representation in Multi-System Regulation} / related use & N058 & V \\
		
		N060 & Representation construction map $F(\Sigma_t)$ & O & L4 & \emph{Emergence of Coherence Representation} & same paper & N016, N023, N024, N053 & S \\
		
		N061 & Downstream regulatory map $\Psi$ & O & L4 & \emph{Emergence of Coherence Representation} & same paper & N058 & V \\
		
		N062 & Regime selection map $\rho$ & O & L4/L5 & \emph{Emergence of Coherence Representation} / \emph{Coherence Representation in Multi-System Regulation} & latter paper for inter-system use & N058 & V \\
		
		N070 & Trajectory-dependent regulation & D & L5 & \emph{Trajectory-Dependent Regulation in Cognitive Systems} & same paper & N025, N026, N028 & S \\
		
		N071 & Identity as regulatory attractor & D & L5 & \emph{Identity as a Regulatory Attractor} & same paper & N024, N027, N070 & S \\
		
		N072 & Identity-core / low-overload attractor region & C & L5 & \emph{Identity as a Regulatory Attractor} & same paper & N071 & S \\
		
		N073 & Inter-system comparability / order alignment & P/C & L5 & \emph{Coherence Representation in Multi-System Regulation} & same paper & N057, N058 & S \\
		
		N074 & Multi-system co-regulation & D/P & L5 & \emph{Coherence Representation in Multi-System Regulation} & same paper & N073, N062 & S \\
		
		N075 & Inter-system conflict geometry & D & L5 & \emph{Inter-System Conflict Geometry in Cognitive Systems} & same paper & N053, N055, N074 & S \\
		
		N076 & Conflict indicator $C_{\mathrm{conf}}(A,B,t)$ & C & L5 & \emph{Inter-System Conflict Geometry in Cognitive Systems} & same paper & N075 & S \\
		
		N077 & General Theory integration node & M & L5 & \emph{General Theory of Cognitive Structuring} & same paper & N013, N024, N026, N031, N050, N057, N039, N071, N075 & G \\
		
		\bottomrule
	\end{longtable}
	
	\newpage
	
	\section{Paper-level dependency table v1}
	
	\begin{longtable}{>{\RaggedRight\arraybackslash}p{0.03\textwidth}
			>{\RaggedRight\arraybackslash}p{0.20\textwidth}
			>{\RaggedRight\arraybackslash}p{0.21\textwidth}
			>{\RaggedRight\arraybackslash}p{0.17\textwidth}
			>{\RaggedRight\arraybackslash}p{0.17\textwidth}
			>{\RaggedRight\arraybackslash}p{0.14\textwidth}
			>{\RaggedRight\arraybackslash}p{0.04\textwidth}}
		\toprule
		\textbf{PID} & \textbf{Paper} & \textbf{Imports} & \textbf{Introduces} & \textbf{Extends} & \textbf{Depends logically on} & \textbf{Status} \\
		\midrule
		\endfirsthead
		\toprule
		\textbf{PID} & \textbf{Paper} & \textbf{Imports} & \textbf{Introduces} & \textbf{Extends} & \textbf{Depends logically on} & \textbf{Status} \\
		\midrule
		\endhead
		
		P01 & \emph{Structural Updating and the Limits of Cognitive Change} &
		adjacent literatures only; no internal formal imports required &
		admissibility of structural updating; structural updating vs processing distinction; availability-of-updating question &
		establishes conceptual entry layer for whole series &
		--- &
		S \\
		
		P02 & \emph{Coherence Evaluation in Cognitive Systems} &
		entry distinction between processing and structural stability pressure is compatible but not formally required &
		configuration space; stability region $U$; coherence; coherence measure as distance; coherence evaluation &
		begins geometric branch of the series &
		P01 optional, but not strictly required &
		S \\
		
		P03 & \emph{Structural Admissibility in Cognitive Systems} &
		processing/updating distinction; architecture state; tension idea &
		admissibility operator $A_t$; regulatory phase space $(T,L)$; stability region in phase space; critical boundary; architectural level &
		formalizes when structural updating becomes admissible &
		P01, and operationally P02 for geometric framing &
		S \\
		
		P04 & \emph{Overload Formation in Cognitive Processing} &
		structural tension; resources; rigidity; processing vs updating distinction &
		compensability threshold; instantaneous overload $u_t$; overload memory $L_t$; path dependence of overload accumulation &
		deepens regulatory layer below admissibility decision &
		P03 in series logic, though parts can be read alongside it &
		S \\
		
		P05 & \emph{Trajectory-Dependent Regulation in Cognitive Systems} &
		tension; overload memory; admissibility operator; stability region in $(T,L)$ &
		trajectory-dependent regulation; explicit regulatory phase-space trajectories; hysteresis and boundary deformation dynamics &
		dynamical analysis of regulatory movement before updating &
		P03, P04 &
		S \\
		
		P06 & \emph{Identity as a Regulatory Attractor} &
		overload memory; admissibility boundary; trajectory dependence; bounded regulation &
		identity as regulatory attractor; identity-core; low-overload attractor regions &
		derives identity from long-run regulatory concentration rather than representation &
		P03, P04, P05 &
		S \\
		
		P07 & \emph{Invariants in Cognitive Architectures} &
		earlier papers used stable structural constraints operationally &
		invariant; invariant structure; interaction graph; local admissibility profiles; pairwise incompatibility; invariant-induced stability geometry &
		retroactively grounds the architectural source of stability regions and conflict &
		P02, P03 &
		S \\
		
		P08 & \emph{Structural Compression in Cognitive Architectures} &
		admissibility of updating; overload criterion; invariants; level notion &
		structural compression; structural simplification distinction; compressed invariant $F(S)$; compression-induced geometry change &
		links invariant transformation to long-run overload reduction and level transition &
		P03, P04, P07 &
		S \\
		
		P09 & \emph{Emergence of Coherence Representation} &
		geometric coherence; overload and trajectory dependence; compression; partial observability of stability geometry &
		coherence representation; expressed regulatory representation; representation-construction map $F(\Sigma_t)$; coherence / evaluation / representation distinction &
		adds representation layer within bounded regulation &
		P02, P04, P05, P08 &
		S \\
		
		P10 & \emph{Coherence Representation in Multi-System Regulation} &
		coherence representation; partial observability; regulatory order induced by representation &
		ordinal adequacy; regulatory preorder; inter-system order alignment; multi-system co-regulation &
		extends representation beyond a single architecture without shared metric &
		P09 &
		S \\
		
		P11 & \emph{Pre-Symbolic Admissibility in Cognitive Systems} &
		coherence representation; structural admissibility; invariants; overload context &
		pre-symbolic signal domain $Z_t$; pre-symbolic admissibility operator $\Pi_t^{ps}$; admissible discrepancy domain $D_t$; separation of regulatory layers &
		inserts an earlier filtering layer before representation &
		P03, P07, P09 &
		S \\
		
		P12 & \emph{Restricted Accessibility of Coherence in Cognitive Systems} &
		geometric coherence; pre-symbolic admissibility; coherence representation &
		restricted accessibility of coherence; non-injectivity of regulatory accessibility; structural asymmetry between global state and regulatory access &
		makes explicit that global coherence need not become regulatorily accessible &
		P02, P09, P11 &
		S \\
		
		P13 & \emph{Inter-System Conflict Geometry in Cognitive Systems} &
		admissible discrepancy domains; restricted accessibility; multi-system setting &
		conflict as incompatibility of admissible discrepancy domains; conflict indicator; compatibility via shared admissible discrepancy structure &
		extends access-layer logic into inter-system incompatibility &
		P10, P11, P12 &
		S \\
		
		P14 & \emph{General Theory of Cognitive Structuring} &
		all major earlier layers: structural updating, coherence, regulatory load/overload, invariants, layered admissibility, compression, representation, identity, inter-system consequences &
		synthetic integration of layered admissibility architecture &
		unifies the series; not a foundational starting node &
		P02--P13 strongly; P01 conceptually &
		G \\
		
		\bottomrule
	\end{longtable}
	
	\section{Minimal unfolding order}
	
	The present tables support the following minimal unfolding order of the series:
	
	\begin{enumerate}[leftmargin=2em]
		\item conceptual distinction between processing and structural updating;
		\item geometric formalization of coherence;
		\item formalization of structural admissibility in the joint regulatory state;
		\item overload formation and trajectory dependence;
		\item formalization of invariants as architectural constraints;
		\item structural compression as admissible invariant transformation;
		\item coherence representation as an internal regulatory variable;
		\item pre-symbolic admissibility as an earlier selective layer;
		\item restricted accessibility of coherence as a consequence of geometric coherence plus admissibility-constrained access;
		\item multi-system representation and inter-system incompatibility;
		\item synthetic integration in the \emph{General Theory of Cognitive Structuring}.
	\end{enumerate}
	
	This unfolding order is sufficient for paper-level dependency reconstruction and is consistent with
	the acyclicity criteria stated in the accompanying public note.
	
	\section{Registry interpretation notes}
	
	\subsection{Stable nodes}
	
	Nodes marked \textbf{S} are considered stable enough for public dependency use. This status is
	assigned where the role of the node is already clearly separated in the current corpus and can
	be used directly in graph construction.
	
	\subsection{Verification nodes}
	
	Nodes marked \textbf{V} are boundary cases where the current corpus distinguishes between:
	\begin{itemize}[leftmargin=2em]
		\item early operational usage,
		\item later explicit formalization,
		\item or multiple nearby formulations that should be kept distinct in future graph refinement.
	\end{itemize}
	
	Such nodes are still useful for the present registry, but their exact boundaries may later be
	sharpened.
	
	\subsection{Graph-organizing nodes}
	
	Nodes marked \textbf{G} are included because they improve readability of the series structure.
	They are useful especially for:
	\begin{itemize}[leftmargin=2em]
		\item layer-level graphs,
		\item synthetic convergence maps,
		\item and paper-level DAG summaries.
	\end{itemize}
	
	They should not automatically be treated as theorem-grade units.
	
	\section{Recommended use of this document}
	
	This document is designed to serve as the immediate companion to the public acyclicity note.
	
	It may be used:
	\begin{itemize}[leftmargin=2em]
		\item as an appendix-level dependency registry;
		\item as the source table for a later node-level DAG;
		\item as the source table for a later paper-level DAG;
		\item as a public audit trail showing where the current dependency reconstruction remains fully stable and where it remains open to minor refinement.
	\end{itemize}
	
	\section{Possible later extensions}
	
	

	
\end{document}
